\setcounter{section}{0}

\Needspace{5\baselineskip}
\hypertarget{kv-cache}{%
\section{KV Cache}\label{kv-cache}}

\Needspace{5\baselineskip}
\hypertarget{i.-the-core-principle-of-kv-cache}{%
\subsection{I. The Core Principle of KV Cache}\label{i.-the-core-principle-of-kv-cache}}

In full attention, generating each token has time and space complexity O(n\^{}2) (using a KV cache during inference reduces the time complexity for generating a single token to O(n), gives O(n\^{}2) time complexity for generating a sequence, and O(n) memory complexity). These costs are the main obstacles to processing long sequences (long contexts). In particular, during inference (not training), generating each token incurs an O(L\^{}2) computational cost. However, much of the computation of the K and V matrices from the input sequence is repeated.

\Needspace{5\baselineskip}
To address the computational time complexity during inference, LLMs commonly use a KV cache:

\Needspace{4\baselineskip}
\begin{enumerate}
\def\labelenumi{\arabic{enumi}.}
\tightlist
\item
  When calculating Q, K, and V from the input
\end{enumerate}

\Needspace{5\baselineskip}
In traditional attention:

\[
Q = XW_Q,\quad K = XW_K,\quad V = XW_V
\]

With a KV cache, when generating token t+1, we only need token t as Q. For K and V, rows 1 through t-1 (the K and V values of tokens 1 through t-1) have already been calculated during previous generation steps and can be stored in memory for direct reuse. Only the K and V values of token t need to be calculated:

\[
q_t = x_tW_Q,\quad k_t = x_tW_K,\quad v_t = x_tW_V
\]

Here, k\_t and v\_t are appended below the previously calculated K and V matrices in memory for subsequent use. This reduces the computational complexity from O(L) to O(1).

\Needspace{4\baselineskip}
\begin{enumerate}
\def\labelenumi{\arabic{enumi}.}
\setcounter{enumi}{1}
\tightlist
\item
  When calculating attention weights
\end{enumerate}

\Needspace{5\baselineskip}
In traditional attention:

\[
\mathrm{Scores} = Q \cdot K^T
\]

\Needspace{5\baselineskip}
Description of the operation:

\begin{itemize}
\tightlist
\item
  The dimensions of \(Q\) are \(t \times d_k\).
\item
  The dimensions of \(K\) are \(t \times d_k\), so those of \(K^T\) are \(d_k \times t\).
\item
  This is a standard matrix multiplication, yielding an attention score matrix of dimensions \(t \times t\).
\end{itemize}

\Needspace{5\baselineskip}
With a KV cache:

\[
\mathrm{Scores} = q_t \cdot (K_{\mathrm{cache}}^{(t)})^T
\]

\Needspace{5\baselineskip}
Description of the operation:

\begin{itemize}
\tightlist
\item
  The dimensions of \(q_t\) are \(1 \times d_k\).
\item
  The dimensions of \(K_{\mathrm{cache}}^{(t)}\) are \(t \times d_k\) (\(d_k \times t\) after transposition).
\item
  This is a ``vector--matrix'' multiplication, yielding attention scores of dimensions \(1 \times t\) for the current step.
\end{itemize}

The computational complexity is therefore reduced from O(L\^{}2) to O(L).

\Needspace{4\baselineskip}
\begin{enumerate}
\def\labelenumi{\arabic{enumi}.}
\setcounter{enumi}{2}
\tightlist
\item
  When calculating the output by weighting
\end{enumerate}

\Needspace{5\baselineskip}
In traditional attention:

\[
Y = \mathrm{Softmax}(\mathrm{Scores}) \cdot V
\]

\Needspace{5\baselineskip}
Description of the operation:

\begin{itemize}
\tightlist
\item
  The attention matrix has dimensions \(t \times t\).
\item
  The \(V\) matrix has dimensions \(t \times d_v\).
\item
  The output \(Y\) has dimensions \(t \times d_v\).
\end{itemize}

\Needspace{5\baselineskip}
With a KV cache:

\[
y_t = \alpha_t \cdot V_{\mathrm{cache}}^{(t)}
\]

\Needspace{5\baselineskip}
Description of the operation:

\begin{itemize}
\tightlist
\item
  The weights \(\alpha_t\) have dimensions \(1 \times t\).
\item
  \(V_{\mathrm{cache}}^{(t)}\) has dimensions \(t \times d_v\).
\item
  This is a ``vector--matrix'' multiplication, and the output \(y_t\) has dimensions \(1 \times d_v\).
\end{itemize}

\Needspace{5\baselineskip}
\hypertarget{ii.-why-is-kv-cache-not-used-during-training}{%
\subsection{II. Why Is KV Cache Not Used during Training?}\label{ii.-why-is-kv-cache-not-used-during-training}}

A KV cache is designed to address the inefficiency of ``serial computation'' when generating one token at a time during inference. Training, however, is naturally parallel: we already know the complete answer, and teacher forcing and masking allow us to calculate the losses at all positions in one computation. There is therefore no need to ``cache the previous step's result for the next step.'' Forcibly using a KV cache during training would instead turn efficient parallel training into inefficient serial training, making training hundreds or thousands of times slower.

As discussed earlier, during pretraining or fine-tuning, the steps are completed simultaneously in a single matrix operation.

We obtain Q, K, and V for all tokens at once, then calculate attention scores and apply a mask. The mask is an ``upper-triangular matrix'' that acts after attention score calculation and before Softmax. It forces the attention scores for all future positions (i,j) (i\textgreater j), that is, the dot products of q\_i and k\_j, to a negative number with an extremely large absolute value (for example, -1e9).

There is no situation in this process where ``step 1 finishes and stores its result for step 2,'' because steps 1 and 2 occur simultaneously.

{[}Special case: during PPO training in RLHF, the procedure has two steps. The rollout generation stage lets the model generate responses to prompts; this is essentially inference, so a KV cache is used. The update stage uses the generated responses and rewards to calculate gradients and update the model, returning to parallel operation without a KV cache.{]}

\Needspace{5\baselineskip}
\hypertarget{iii.-prefill}{%
\subsection{III. Prefill}\label{iii.-prefill}}

Before generation starts, we let the model read the entire prompt in parallel in one pass, calculate the initial Q, K, and V matrices using Q = X\emph{W\_Q, K = X}W\_K, and V = X*W\_V, and store the initial KV cache. Decoding then begins, generating tokens autoregressively.

Prefill is computation-intensive, whereas decoding is memory-intensive. Some architectures assign the two stages to GPUs with high computational power and high GPU-memory bandwidth, respectively.

\Needspace{5\baselineskip}
\hypertarget{iv.-system-hardware-management-paged-attention-vllm}{%
\subsection{IV. System Hardware Management: Paged Attention (vLLM)}\label{iv.-system-hardware-management-paged-attention-vllm}}

A KV cache requires contiguous space in GPU memory, causing considerable waste from memory fragmentation (similar to fragmentation in operating-system memory). Borrowing the paging technique from operating-system virtual memory, it stores the KV cache in blocks in noncontiguous physical GPU memory and indexes them through a page table. This greatly improves GPU-memory utilization, allowing a larger batch size with the same memory and thereby improving inference throughput.

\FloatBarrier
\Needspace{5\baselineskip}
\hypertarget{references-91}{%
\subsection{References}\label{references-91}}

\vspace{0.25em}
\begin{itemize}
\tightlist
\item
  Kwon, W., Li, Z., Zhuang, S., et al.~(2023). \href{https://arxiv.org/abs/2309.06180}{Efficient Memory Management for Large Language Model Serving with PagedAttention}. SOSP 2023.
\end{itemize}

\clearpage
